% =============================================================================
%  theorem_49_completion.tex
%
%  A COMPLETE PROOF of Theorem 49
%  (\Cref{theorem:ramification_after_blowup_is_tame_for_p_power}):
%
%      For G = Z/p^r Z, the symmetric-power torsor P^{[deg n]} is
%      unramified at the generic point of the exceptional divisor of the
%      blowup of C^{(deg n)} at n = nP.
%
%  The proof is direct (no induction on r, no Swan conductors).  It replaces
%  the single Artin--Schreier equation of the thesis's Z/p computation
%  (RamificationAfterBlowupIsTame.tex, lines 499-575) by a length-r Witt
%  vector, and consists of three steps:
%
%    (1)  The Artin--Schreier--Witt datum of P^{[d]} at the exceptional
%         divisor is the d-fold WITT SUM of the local datum of P
%         (Theorem \ref{thm:datum} below -- proved by an explicit,
%         elementary change of variables; this closes open point (1) of
%         q_ramification_along_VF/general_witt_proof/general_witt_proof.tex).
%
%    (2)  The Witt-sum polynomials are ISOBARIC, so each component of the
%         datum is a symmetric polynomial in y_i = 1/x_i of total degree
%         at most the largest upper ramification break u < d
%         (Lemma \ref{lem:isobaric-bound}).
%
%    (3)  A symmetric polynomial in the y_i of degree < d is a polynomial
%         in e_1(y),...,e_{d-1}(y) = e_{d-1}(x)/e_d(x),...,e_1(x)/e_d(x),
%         hence lies in the RESIDUE FIELD of the exceptional valuation --
%         so the whole datum is integral, and the extension is unramified
%         because the Artin--Schreier--Witt isogeny is etale
%         (Lemma \ref{lem:integrality}, Corollary \ref{cor:integral-unramified}).
%
%  The only non-elementary input is the classical Schmid--Witt description
%  of the ramification breaks of Artin--Schreier--Witt extensions
%  (\cite{Thomas2005}), which the thesis already cites for r = 1
%  (theorem:artin_schreier_ramification).  Everything else is proved here.
%
%  Sections 1-2 (Witt vectors, ASW theory) are background intended for the
%  Preliminaries; Section 3-4 are the proof itself, intended to replace the
%  broken inductive proof in RamificationAfterBlowupIsTame.tex (lines
%  636-673).  See the final remarks for integration notes.
%
%  Compile: TEXINPUTS=../..: latexmk -pdf theorem_49_completion.tex  (biber)
% =============================================================================
\documentclass[11pt]{article}
\usepackage{geometry}\geometry{margin=1in}
\usepackage[T1]{fontenc}
\usepackage{ifthen}

\setlength{\parindent}{1em}
\setlength{\parskip}{0.5em}

\input{includes}
\input{MacrosAndFigures}

\usepackage{csquotes}
\usepackage[
backend=biber,
style=alphabetic,
sorting=ynt
]{biblatex}
\addbibresource{Bib.bib}

\newboolean{ThesisIntro}
\setboolean{ThesisIntro}{false}

% ---- local notation ---------------------------------------------------------
\providecommand{\eps}{\varepsilon}                  % elementary symmetric in y
\newcommand{\wv}[1]{\underline{#1}}                 % a Witt vector
\newcommand{\Wsum}{\sideset{}{^{W}}\sum}            % Witt-vector sum
\newcommand{\mm}{\mathfrak{m}_{k[[x]]}}             % maximal ideal x k[[x]]

\title{Completion of the proof of
\texorpdfstring{\Cref{theorem:ramification_after_blowup_is_tame_for_p_power}}{Theorem 49}:\\
ramification after blowup is trivial for $G=\Z/p^r\Z$\\[4pt]
{\large a direct proof via Artin--Schreier--Witt theory}}
\author{}
\date{\today}

\begin{document}
\maketitle

\noindent\textbf{Goal.}
Let $C$ be a projective smooth geometrically connected curve over $k$
($\ch k = p$), $\mdls=\sum n_iP_i$ a modulus, $U=C\setminus\mdls$, and let
$\cP\to U$ be a $G$-torsor with ramification bounded by $\mdls$, where
$G=\Z/p^r\Z$.  Fix $P=P_i$, put $n=n_i$, $\ndls=nP$ and $d=\deg\ndls=n$.  Let
$X_\ndls\to C^{(d)}$ be the blowup at the point $\ndls\in C^{(d)}$, with
exceptional divisor $E_\ndls$ and generic point $\eta_\ndls$.  We prove:

\begin{theorem*}[{= \Cref{theorem:ramification_after_blowup_is_tame_for_p_power}}]
$\cP^{[d]}$ is unramified at $\eta_\ndls$.
\end{theorem*}

\noindent\textbf{Strategy.}
The thesis already proves this for $r=1$ by an explicit computation: the
Artin--Schreier datum of $\cP^{[d]}$ at $\eta_\ndls$ is the power sum
$\sum_{i=1}^d f(x_i)$, which is a symmetric function of low degree in the
$y_i=1/x_i$ and therefore lands in the \emph{residue field} of the
exceptional valuation.  The inductive route from $\Z/p$ to $\Z/p^r$ stalls
(see \texttt{a\_wrong\_path.tex}); instead we redo the $r=1$ computation
verbatim with a length-$r$ \emph{Witt vector} in place of a single
Artin--Schreier equation.  The datum of $\cP^{[d]}$ becomes the $d$-fold
\emph{Witt sum} of the local datum of $\cP$; the Witt carries are controlled
by the isobaric grading of the Witt addition polynomials, and the
ramification bound $<d$ makes every component of the datum a symmetric
polynomial of degree $<d$ in the $y_i$ --- hence integral at $\eta_\ndls$,
hence unramified, because the Artin--Schreier--Witt isogeny is \'etale.

No Swan conductors are used.  The only imported result is the classical
Schmid--Witt computation of ramification breaks of Artin--Schreier--Witt
extensions (\cite{Thomas2005}), whose $r=1$ case is
\Cref{theorem:artin_schreier_ramification} of the thesis.
\Cref{sec:witt,sec:asw} are supplementary background (intended for the
Preliminaries); \Cref{sec:datum,sec:proof} contain the proof.

% #############################################################################
\section{Background I: Witt vectors}\label{sec:witt}
% #############################################################################

Fix a prime $p$.  We recall the $p$-typical Witt vectors of finite length
$r\ge 1$; the standard reference is \cite[II \S6]{serreLF} (see also
\cite{Rabinoff2014Witt}).

For a commutative ring $A$, the set $W_r(A)=A^r$ of vectors
$\wv a=(a_0,\dots,a_{r-1})$ carries a commutative ring structure, functorial
in $A$, uniquely determined by the requirement that the \emph{ghost maps}
\[
w_n\colon W_r(A)\longrightarrow A,\qquad
w_n(\wv a)=\sum_{l=0}^{n}p^{\,l}a_l^{\,p^{\,n-l}}
=a_0^{p^n}+pa_1^{p^{n-1}}+\dots+p^na_n
\qquad(0\le n\le r-1)
\]
are ring homomorphisms.  Concretely, there are universal polynomials with
\emph{integer} coefficients computing the ring operations.  We will only
need addition and subtraction, and only their following structural features.

\begin{notation*}
For $d\ge2$ let
$S^{(d)}_n\in\Z\big[X^{(i)}_l : 1\le i\le d,\ 0\le l\le r-1\big]$,
$0\le n\le r-1$, denote the $d$-fold Witt \emph{addition} polynomials,
characterized by
\begin{equation}\label{eq:ghost-sum}
w_n\big(S^{(d)}_0,\dots,S^{(d)}_{n}\big)
=\sum_{i=1}^{d}w_n\big(\wv X^{(i)}\big),
\end{equation}
so that
$\wv X^{(1)}+_W\dots+_W\wv X^{(d)}
=\big(S^{(d)}_0,\dots,S^{(d)}_{r-1}\big)$.
Similarly let $D_n(\wv X;\wv Y)$ denote the subtraction polynomials,
characterized by $w_n(\wv D)=w_n(\wv X)-w_n(\wv Y)$.
We write $+_W$, $-_W$ and $\Wsum$ for Witt-vector operations, to
distinguish them from componentwise ones, and $\wv 0=(0,\dots,0)$.
\end{notation*}

Expanding \cref{eq:ghost-sum} and using
$w_n(\wv X)=p^nX_n+\sum_{l<n}p^lX_l^{p^{n-l}}$ we get the recursion
(valid in $\Q[X^{(i)}_l]$, where the polynomials are uniquely determined):
\begin{equation}\label{eq:recursion}
p^{\,n}S^{(d)}_n
=\sum_{i=1}^{d}w_n\big(\wv X^{(i)}\big)
-\sum_{l=0}^{n-1}p^{\,l}\big(S^{(d)}_l\big)^{p^{\,n-l}},
\end{equation}
and the analogous recursion for $D_n$ with
$\sum_i w_n(\wv X^{(i)})$ replaced by $w_n(\wv X)-w_n(\wv Y)$.
That these recursions produce polynomials with integer coefficients is
Witt's theorem \cite[II \S6, Thm.~6]{serreLF}.

\begin{definition}[weights]\label{def:weight}
Assign to the level-$l$ variable $X^{(i)}_l$ (resp.\ $X_l$, $Y_l$) the
\emph{weight} $p^{\,l}$.  The weight of a monomial is the sum of the weights
of its variables, counted with multiplicity.  A polynomial is
\emph{isobaric of weight $w$} if every monomial occurring in it has weight
exactly $w$.  (In particular an isobaric polynomial of weight $w>0$ has no
constant term.)
\end{definition}

\begin{lemma}[isobaricity]\label{lem:isobaric}
$S^{(d)}_n$ and $D_n$ are isobaric of weight $p^{\,n}$.
\end{lemma}

\begin{proof}
The ghost polynomial $w_n$ is isobaric of weight $p^n$: its monomial
$p^lX_l^{\,p^{n-l}}$ has weight $p^{\,l}\cdot p^{\,n-l}=p^{\,n}$.  We induct
on $n$ using \cref{eq:recursion}.  For $n=0$,
$S^{(d)}_0=\sum_iX^{(i)}_0$ is isobaric of weight $1$.  For $n>0$: each
$w_n(\wv X^{(i)})$ is isobaric of weight $p^n$; by induction $S^{(d)}_l$ is
isobaric of weight $p^{\,l}$ for $l<n$, hence $(S^{(d)}_l)^{p^{n-l}}$ is
isobaric of weight $p^{\,l}\cdot p^{\,n-l}=p^{\,n}$.  So the right-hand side
of \cref{eq:recursion} is isobaric of weight $p^n$ in $\Q[X^{(i)}_l]$, and
dividing by $p^{\,n}$ does not change the set of monomials.  The same
induction applies to $D_n$.
\end{proof}

\begin{lemma}[triangularity]\label{lem:triangular}
\begin{enumerate}
  \item $S^{(d)}_n=\sum_{i=1}^{d}X^{(i)}_n+T_n$ with
        $T_n\in\Z\big[X^{(i)}_l : l<n\big]$, and
        $D_n=X_n-Y_n+E_n$ with $E_n\in\Z[X_l,Y_l : l<n]$.
        In words: \emph{the $n$-th component of a Witt sum or difference is
        the plain sum or difference of the $n$-th components, plus an
        integer polynomial in the components of levels $<n$.}
  \item Consequently, if $\wv u,\wv v\in W_r(A)$ satisfy $u_l=v_l=0$ for all
        $l<n$, then $\wv u\pm_W\wv v$ agrees with the componentwise sum or
        difference up to and including level $n$.
\end{enumerate}
\end{lemma}

\begin{proof}
(1) Induct on $n$ via \cref{eq:recursion}: in
$w_n(\wv X^{(i)})=p^nX^{(i)}_n+(\text{polynomial in }X^{(i)}_{l},\,l<n)$
only the displayed term involves level $n$, and the correction terms
$(S^{(d)}_l)^{p^{n-l}}$, $l<n$, involve only levels $\le l<n$ by induction.
Divide by $p^n$; integrality of the result is part of Witt's theorem.
(2) By (1), the components of level $\le n$ of $\wv u\pm_W \wv v$ are
$u_l\pm v_l+(\text{polynomial without constant term in levels}<l)$,
and the correction vanishes when all lower components vanish
(\Cref{lem:isobaric} guarantees there is no constant term).
\end{proof}

\subsubsection*{Frobenius and the operator \texorpdfstring{$\wp$}{p}}

From now on $A$ is an $\F_p$-algebra.  The absolute Frobenius
$\phi\colon A\to A$, $a\mapsto a^p$, is a ring endomorphism, so by
functoriality it induces a ring endomorphism
\[
F:=W_r(\phi)\colon W_r(A)\to W_r(A),\qquad
F(a_0,\dots,a_{r-1})=(a_0^{\,p},\dots,a_{r-1}^{\,p}),
\]
which acts \emph{componentwise}.  (For $\F_p$-algebras this componentwise
map coincides with the Witt-vector Frobenius of the general theory
\cite[II \S6]{serreLF}; we will only ever use the componentwise
description, which we may take as the definition.)  Define the
\emph{Artin--Schreier--Witt operator}
\[
\wp:=F-\mathrm{id}\colon W_r(A)\longrightarrow W_r(A),
\qquad \wp(\wv a)=F(\wv a)-_W\wv a .
\]
Since $F$ is a ring homomorphism, $\wp$ is an endomorphism of the additive
group $W_r(A)$, functorial in $A$.  For $r=1$ it is the classical
Artin--Schreier operator $a\mapsto a^p-a$.

\begin{lemma}[shape of $\wp$]\label{lem:wp-triangular}
For variables $\wv Y=(Y_0,\dots,Y_{r-1})$ over $\F_p$,
\[
\wp(\wv Y)_n \;=\; Y_n^{\,p}-Y_n+\Delta_n(Y_0,\dots,Y_{n-1}),
\qquad 0\le n\le r-1,
\]
where $\Delta_n\in\F_p[Y_0,\dots,Y_{n-1}]$ has no constant term and
$\Delta_0=0$.
\end{lemma}

\begin{proof}
$\wp(\wv Y)_n=D_n\big(Y_0^{\,p},\dots,Y_{r-1}^{\,p};\,Y_0,\dots,Y_{r-1}\big)$.
By \Cref{lem:triangular}(1) this equals
$Y_n^{\,p}-Y_n+E_n(Y^{\,p}_{<n};Y_{<n})$, and we set
$\Delta_n:=E_n(Y^{\,p}_{<n};Y_{<n})\bmod p$, a polynomial in
$Y_0,\dots,Y_{n-1}$ only.  It has no constant term because $E_n$ has none
(\Cref{lem:isobaric}).
\end{proof}

\begin{lemma}[kernel of $\wp$]\label{lem:kernel}
If $A$ is an integral domain over $\F_p$, then
$\ker\big(\wp\colon W_r(A)\to W_r(A)\big)=W_r(\F_p)\cong\Z/p^r\Z$.
\end{lemma}

\begin{proof}
$\wp(\wv a)=\wv 0$ iff $F\wv a=\wv a$, i.e.\ $a_l^{\,p}=a_l$ for all $l$
(both maps are componentwise), i.e.\ $a_l\in\F_p$ for all $l$ since $A$ is a
domain.  Finally $W_r(\F_p)=W(\F_p)/V^rW(\F_p)=\Z_p/p^r\Z_p\cong\Z/p^r\Z$
\cite[II \S6]{serreLF}.
\end{proof}

From now on we identify $G:=\Z/p^r\Z=W_r(\F_p)$.

\begin{prop}[the Artin--Schreier--Witt isogeny is \'etale]\label{prop:etale}
Let $A$ be an $\F_p$-algebra and $\wv a\in W_r(A)$.  Set
\[
A_{\wv a}\;:=\;A[Y_0,\dots,Y_{r-1}]\,/\,
\big(\,\wp(\wv Y)_n-a_n\ :\ 0\le n\le r-1\,\big).
\]
Then:
\begin{enumerate}
  \item $A_{\wv a}$ is a free $A$-module of rank $p^{\,r}$, with basis the
        monomials $\prod_l Y_l^{\,c_l}$, $0\le c_l\le p-1$;
  \item $A_{\wv a}$ is \'etale over $A$;
  \item the rule $\wv Y\mapsto\wv Y+_W\wv g$, $\wv g\in G=W_r(\F_p)$,
        defines an action of $G$ on $A_{\wv a}$ by $A$-algebra
        automorphisms, and if $A\ne 0$ this makes
        $\cP_{\wv a}:=\Spec A_{\wv a}$ a $G$-torsor over $\Spec A$ in the
        \'etale topology.
\end{enumerate}
\end{prop}

\begin{proof}
By \Cref{lem:wp-triangular} the defining equations are
\[
Y_n^{\,p}-Y_n \;=\; a_n-\Delta_n(Y_0,\dots,Y_{n-1}),\qquad n=0,\dots,r-1,
\]
a \emph{triangular tower} of Artin--Schreier equations: each is monic of
degree $p$ in $Y_n$ with coefficients in
$A[Y_0,\dots,Y_{n-1}]$.  (1) follows at once.

(2) $A_{\wv a}$ is flat (free) over $A$; it remains to see it is
unramified, i.e.\ $\Omega^1_{A_{\wv a}/A}=0$.  The module
$\Omega^1_{A_{\wv a}/A}$ is generated by $dY_0,\dots,dY_{r-1}$ subject to
the relations $\sum_m\big(\partial g_n/\partial Y_m\big)\,dY_m=0$, where
$g_n=Y_n^p-Y_n+\Delta_n-a_n$.  The Jacobian matrix
$\big(\partial g_n/\partial Y_m\big)_{n,m}$ is lower triangular with
diagonal entries $pY_n^{\,p-1}-1=-1$, hence invertible; so
$\Omega^1_{A_{\wv a}/A}=0$.

(3) For $\wv g\in W_r(\F_p)$ we have $F\wv g=\wv g$, so
$\wp(\wv Y+_W\wv g)=\wp(\wv Y)+_W\wp(\wv g)=\wp(\wv Y)$; hence the
substitution $\wv Y\mapsto\wv Y+_W\wv g$ preserves the defining ideal
and gives an $A$-algebra endomorphism, with inverse given by $-\wv g$; and
$(\wv g,\wv g')\mapsto\wv g+_W\wv g'$ makes this an action of $G$.
For the torsor property it suffices, since $\cP_{\wv a}\to\Spec A$ is
finite \'etale and surjective of degree $p^r=|G|$, to check that $G$ acts
simply transitively on geometric fibers.  Let $\Omega$ be an algebraically
closed $A$-field; the $\Omega$-points of the fiber are the solutions
$\wv b\in W_r(\Omega)$ of $\wp(\wv b)=\wv a$.  If $\wv b,\wv b'$ are two
solutions then $\wp(\wv b'-_W\wv b)=\wv 0$, so
$\wv b'-_W\wv b\in W_r(\F_p)=G$ by \Cref{lem:kernel}: the action is
transitive; and it is free because translation by $\wv g\ne\wv 0$ in the
group $W_r(\Omega)$ has no fixed points.
\end{proof}

\begin{cor}[integral datum $\Rightarrow$ unramified]\label{cor:integral-unramified}
Let $\Oo$ be a discrete valuation ring containing $\F_p$, with fraction
field $K$.  If $\wv a\in W_r(\Oo)\subset W_r(K)$, then every field factor
of the \'etale $K$-algebra $K_{\wv a}=K\otimes_\Oo\Oo_{\wv a}$ is an
\emph{unramified} extension of $K$ with respect to $\Oo$ (in the sense of
\Cref{definition:tamely_ramified_dvrs}).
\end{cor}

\begin{proof}
By \Cref{prop:etale}, $B:=\Oo_{\wv a}$ is finite \'etale over $\Oo$.  Being
\'etale over a normal domain, $B$ is a finite product $\prod_j B_j$ of
normal domains, each finite \'etale over $\Oo$; then
$B_j$ is the integral closure of $\Oo$ in the field factor
$L_j=\Frac(B_j)$ of $K_{\wv a}$, it is a semi-local Dedekind domain, and at
each of its maximal ideals the extension of discrete valuation rings is
flat and unramified: the ramification index is $1$ and the residue field
extension is finite separable.  This is precisely the definition of
$L_j/K$ being unramified with respect to $\Oo$.
\end{proof}

% #############################################################################
\section{Background II: Artin--Schreier--Witt theory}\label{sec:asw}
% #############################################################################

\subsection{Classification of \texorpdfstring{$\Z/p^r\Z$}{Z/p\^{}r}-torsors on affine schemes}

\begin{prop}\label{prop:classification}
Let $U=\Spec A$ be an integral affine scheme over $\F_p$.  Then
\Cref{prop:etale} induces a bijection
\[
W_r(A)\,/\,\wp\,W_r(A)\;\xrightarrow{\ \cong\ }\;
\mathrm{Tors}\big(U_{\text{\'et}},\Z/p^r\Z\big),\qquad
[\wv a]\longmapsto\big[\cP_{\wv a}\big].
\]
Moreover the bijection is functorial: for a morphism
$\Spec B\to\Spec A$ the pullback of $\cP_{\wv a}$ is
$\cP_{\wv a'}$ where $\wv a'$ is the image of $\wv a$ in $W_r(B)$.
\end{prop}

\begin{proof}
The sequence of abelian sheaves on $U_{\text{\'et}}$
\[
0\longrightarrow \Z/p^r\Z\longrightarrow W_r(\Oo_U)
\xrightarrow{\ \wp\ } W_r(\Oo_U)\longrightarrow 0
\]
is exact: the kernel is the constant sheaf $\Z/p^r\Z$ by
\Cref{lem:kernel} (sections over a connected affine \'etale $V\to U$ with
$V=\Spec B$, $B$ a domain --- and every \'etale $V$ is a disjoint union of
such); and $\wp$ is surjective as a map of \'etale sheaves because for any
$\wv a\in W_r(B)$ the finite \'etale surjective cover
$\Spec B_{\wv a}\to\Spec B$ of \Cref{prop:etale} carries a tautological
solution $\wv Y$ of $\wp(\wv Y)=\wv a$.

The long exact cohomology sequence gives
\[
W_r(A)\xrightarrow{\ \wp\ }W_r(A)\xrightarrow{\ \delta\ }
H^1\big(U_{\text{\'et}},\Z/p^r\Z\big)\longrightarrow
H^1\big(U_{\text{\'et}},W_r(\Oo_U)\big),
\]
and the boundary map $\delta$ sends $\wv a$ to the class of the sheaf of
solutions of $\wp(\wv Y)=\wv a$, i.e.\ of the torsor $\cP_{\wv a}$.
It remains to show $H^1(U_{\text{\'et}},W_r(\Oo_U))=0$.  For $r=1$,
$W_1(\Oo_U)=\bG_a$ and
$H^1(U_{\text{\'et}},\bG_a)\cong H^1(U,\Oo_U)=0$ since $U$ is affine
(\'etale and Zariski cohomology of quasi-coherent sheaves agree,
\cite[III \S3]{milneEC}).  For $r>1$, \Cref{lem:triangular}(2) shows that
$\{(0,\dots,0,a)\}\subset W_r$ is a subgroup isomorphic to $\bG_a$ and that
truncation $(a_0,\dots,a_{r-1})\mapsto(a_0,\dots,a_{r-2})$ is a group
homomorphism (by \Cref{lem:triangular}(1), level-$n$ components of a sum
depend only on levels $\le n$), giving an exact sequence of \'etale sheaves
$0\to\bG_a\to W_r(\Oo_U)\to W_{r-1}(\Oo_U)\to 0$; conclude by induction.
Functoriality of the construction $\wv a\mapsto A_{\wv a}$ is clear.
\end{proof}

\begin{lemma}[torsor invariants]\label{lem:torsor-invariants}
Let $\Gamma$ be a finite group and $\Spec B\to\Spec A$ a $\Gamma$-torsor.
Then $B^{\Gamma}=A$.
\end{lemma}

\begin{proof}
$A\to B$ is finite \'etale, in particular faithfully flat, and the torsor
condition gives a $\Gamma$-equivariant isomorphism of $B$-algebras
$B\otimes_AB\cong\mathrm{Maps}(\Gamma,B)$,
$b\otimes b'\mapsto(\gamma\mapsto\gamma(b)\,b')$, where $\Gamma$ acts on
the left tensor factor and, on the right side, by permuting arguments
($\gamma_0$ sends a map $\varphi$ to $\gamma\mapsto\varphi(\gamma\gamma_0)$).
Taking $\Gamma$-invariants, which commutes with the flat base change
$A\to B$,
\[
B^\Gamma\otimes_AB=(B\otimes_AB)^{\Gamma\otimes 1}
\cong\mathrm{Maps}(\Gamma,B)^{\Gamma}=\{\text{constant maps}\}=B
=A\otimes_AB .
\]
By faithful flatness of $B$ over $A$, $B^\Gamma=A$.
\end{proof}

\begin{lemma}[additivity: contracted products]\label{lem:additivity}
In the situation of \Cref{prop:classification}, for
$\wv a,\wv b\in W_r(A)$ there is a canonical isomorphism of $G$-torsors
\[
\cP_{\wv a}\otimes^{G}\cP_{\wv b}\;\cong\;\cP_{\wv a+_W\wv b}.
\]
\end{lemma}

\begin{proof}
Write the contracted product as
$(\cP_{\wv a}\times_U\cP_{\wv b})/G^{-}$, where
$G^-=\{(\wv g,-\wv g)\}\le G\times G$; the ring of
$\cP_{\wv a}\times_U\cP_{\wv b}$ is
\[
T=A[\wv X_1,\wv X_2]\,/\,
\big(\wp(\wv X_1)-\wv a,\ \wp(\wv X_2)-\wv b\big),
\]
where, here and below, ``$\wp(\wv X)-\wv c$'' stands for the $r$
componentwise equations, and $\wv X_1,\wv X_2$ are two batches of $r$
variables.

\emph{Change of variables.}  Put $\wv Y:=\wv X_1+_W\wv X_2$.  By
\Cref{lem:triangular}(1),
$Y_l=X_{1,l}+\big(\text{polynomial in }X_{1,<l}\text{ and }X_{2,\le l}\big)$,
so the substitution $X_{1,l}\mapsto Y_l$ ($l=0,\dots,r-1$) is a
triangular, unipotent change of polynomial generators:
$A[\wv X_1,\wv X_2]=A[\wv Y,\wv X_2]$.

\emph{The ideal in the new variables.}  We claim
\[
\big(\wp(\wv X_1)-\wv a,\ \wp(\wv X_2)-\wv b\big)
=\big(\wp(\wv Y)-(\wv a+_W\wv b),\ \wp(\wv X_2)-\wv b\big).
\]
Both inclusions follow from the elementary fact that for any polynomial
$P$ and tuples $u,v$ one has $P(u)-P(v)\in(u_1-v_1,u_2-v_2,\dots)$,
combined with Witt-vector identities:
for ``$\subseteq$'' use $\wv X_1=\wv Y-_W\wv X_2$ and the identity
$\wp(\wv Y-_W\wv X_2)=\wp(\wv Y)-_W\wp(\wv X_2)$ (as polynomial identities
in the components --- $\wp$ is additive), evaluating $P=$ each component of
$\wp(\,\cdot\,-_W\,\cdot\,)$ at
$u=(\wp\wv Y,\wp\wv X_2)$ and $v=(\wv a+_W\wv b,\wv b)$, and noting
$(\wv a+_W\wv b)-_W\wv b=\wv a$; the inclusion ``$\supseteq$'' is the same
computation run backwards, with $P=$ components of
$\wp(\,\cdot\,+_W\,\cdot\,)$.  Hence
\[
T\;\cong\;A_{\wv a+_W\wv b}\ \otimes_A\ A_{\wv b},
\qquad \wv Y\mapsto \wv Y\otimes 1,\quad \wv X_2\mapsto 1\otimes\wv X_2 .
\]
\emph{Invariants.}  The subgroup $G^-$ acts by
$\wv Y\mapsto\wv Y+_W\wv g-_W\wv g=\wv Y$ and
$\wv X_2\mapsto\wv X_2-_W\wv g$; i.e.\ trivially on the first tensor
factor and through the full torsor action of $G$ on the second.  By
\Cref{lem:torsor-invariants},
$T^{G^-}=A_{\wv a+_W\wv b}\otimes_A(A_{\wv b})^{G}
=A_{\wv a+_W\wv b}$.
The residual $G=(G\times G)/G^-$-action is
$\wv Y\mapsto\wv Y+_W\wv g$, as required.
\end{proof}

\subsection{Local theory over \texorpdfstring{$k((x))$}{k((x))}: reduced data and ramification breaks}

Let $k$ be an algebraically closed field of characteristic $p$, let
$K=k((x))$ with valuation $v$, valuation ring $k[[x]]$ and maximal ideal
$\mm=xk[[x]]$.  By \Cref{prop:classification} applied to Galois cohomology
(or see \cite{Thomas2005}), continuous characters
$\rho\colon\Gal(K^{sep}/K)\to\Z/p^r\Z$ correspond to classes
$[\wv f]\in W_r(K)/\wp W_r(K)$; the fixed field of $\ker\rho$ is the
splitting field of $\wp(\wv Y)=\wv f$.

\begin{definition}[reduced datum]\label{def:reduced}
A Witt vector $\wv f=(f_0,\dots,f_{r-1})\in W_r(K)$ is \emph{reduced} if
every component $f_l$ is a polynomial in $x^{-1}$ without constant term,
and its degree $m_l:=\deg_{x^{-1}}f_l$ (the pole order of $f_l$; set
$m_l=0$ if $f_l=0$) is either $0$ or prime to $p$.
We call
\[
\mu(\wv f)\;:=\;\max_{0\le l\le r-1}\ p^{\,r-1-l}\,m_l
\]
the \emph{weighted pole order} of $\wv f$.
\end{definition}

\begin{lemma}[reduction to reduced form]\label{lem:reduction}
Every class in $W_r(K)/\wp W_r(K)$ has a reduced representative.
\end{lemma}

\begin{proof}
We first record two solvability facts, both immediate from the triangular
shape of $\wp$ (\Cref{lem:wp-triangular}).

\emph{(a) $\wp\colon W_r(\mm)\to W_r(\mm)$ is surjective.}
Note $W_r(\mm):=\{\wv c: c_l\in\mm\ \forall l\}$ is a subgroup of
$W_r(k[[x]])$, since the addition polynomials have no constant term.
Given $\wv c\in W_r(\mm)$, solve $b_n^{\,p}-b_n=c_n-\Delta_n(b_{<n})$
recursively: assuming $b_0,\dots,b_{n-1}\in\mm$, the right-hand side
$t:=c_n-\Delta_n(b_{<n})$ lies in $\mm$ ($\Delta_n$ has no constant term),
and $b_n:=-\big(t+t^{\,p}+t^{\,p^2}+\dots\big)\in\mm$ converges $x$-adically
and satisfies $b_n^{\,p}-b_n=t$.

\emph{(b) $\wp\colon W_r(k)\to W_r(k)$ is surjective}, since $k$ is
algebraically closed and each equation
$b_n^{\,p}-b_n=c_n-\Delta_n(b_{<n})$ is solvable in $k$.

Now let $\wv a\in W_r(K)$.  We normalize the components from level $0$
upwards; the key point, guaranteed by \Cref{lem:triangular}(2), is that
subtracting $\wp\big(V^n(h)\big)$, where
$V^n(h):=(0,\dots,0,h,0,\dots,0)$ ($h$ in level $n$), changes the level-$n$
component by exactly $-\wp(h)=-(h^p-h)$, changes \emph{no} component of
level $<n$, and changes the components of level $>n$ by elements of $K$.

So suppose levels $0,\dots,n-1$ are already reduced.  Write
$a_n=a_n^-+c+a_n^+$ with $a_n^-\in x^{-1}k[x^{-1}]$, $c\in k$,
$a_n^+\in\mm$.
\begin{itemize}
  \item By (a) and (b) choose $h',h''$ with $\wp(h')=a_n^+$
        ($h'\in\mm$) and $\wp(h'')=c$ ($h''\in k$); subtracting
        $\wp\big(V^n(h'+h'')\big)$ we may assume $a_n=a_n^-$.
  \item If the pole order $m$ of $a_n$ is positive and divisible by $p$,
        with leading term $cx^{-m}$, subtract
        $\wp\big(V^n(c^{1/p}x^{-m/p})\big)$: the level-$n$ component
        changes by $-(cx^{-m}-c^{1/p}x^{-m/p})$, which strictly decreases
        its pole order.  Iterate; after finitely many steps, and after
        removing once more the constant and positive parts as above, the
        level-$n$ component is a polynomial in $x^{-1}$ without constant
        term, of degree $0$ or prime to $p$.
\end{itemize}
All modifications leave levels $<n$ untouched, so the process terminates
after treating level $r-1$.
\end{proof}

The following classical theorem of Schmid and Witt describes the
ramification of the corresponding extensions; it is the promised
Artin--Schreier--Witt input.  For $r=1$ it is exactly
\Cref{theorem:artin_schreier_ramification}.

\begin{theorem}[Schmid--Witt; {\cite{Thomas2005}}]\label{thm:schmid-witt}
Let $\wv f\in W_r(K)$ be reduced with weighted pole order $\mu=\mu(\wv f)$,
and let $\rho\colon\Gal(K^{sep}/K)\to\Z/p^r\Z$ be the character of the
class $[\wv f]$.  Then, in the upper-numbering convention of
\Cref{theorem:artin_schreier_ramification}:
\begin{enumerate}
  \item if $\mu=0$ then $\rho\big(\Gal(K^{sep}/K)^{(u)}\big)=1$ for all
        $u\ge 1$ (the extension cut out by $\rho$ is unramified);
  \item if $\mu>0$ then the largest upper ramification break of the
        extension cut out by $\rho$ equals $\mu$; that is,
        $\rho\big(\Gal(K^{sep}/K)^{(\mu)}\big)\neq 1$ and
        $\rho\big(\Gal(K^{sep}/K)^{(u)}\big)=1$ for $u>\mu$.
\end{enumerate}
In particular, in the sense of \Cref{definition:local_ramification_boundness},
the class $[\wv f]$ has ramification bounded by $d$ if and only if
$\mu(\wv f)\le d-1$.
\end{theorem}

\begin{rem*}
Only the implication ``ramification bounded by $d$ $\Rightarrow$
$\mu\le d-1$'' (i.e.\ the sharpness statement in (2)) is used below; the
easy converse direction is not needed.  Note also that
\Cref{thm:schmid-witt} does not require $\rho$ to be surjective: if the
extension cut out by $\rho$ has degree $p^{r'}<p^r$, the statement is
consistent with itself under the shift identification
$W_{r'}\hookrightarrow W_r$ (Verschiebung), since
$\mu$ is invariant under it.
\end{rem*}

% #############################################################################
\section{The Artin--Schreier--Witt datum of the symmetric power}\label{sec:datum}
% #############################################################################

\subsection{Setting and reductions}\label{subsec:reductions}

We place ourselves in the situation of
\Cref{theorem:ramification_after_blowup_is_tame_for_p_power}, whose
notation we keep: $G=\Z/p^r\Z$; $\cP\to U=C\setminus\mdls$ has ramification
bounded by $\mdls$; $P\in\Supp\mdls$ with multiplicity $n$;
$\ndls=nP$ and $d:=\deg\ndls=n$; and $\eta_\ndls$ is the generic point of
the exceptional divisor of the blowup of $C^{(d)}$ at $\ndls$.

As in the thesis we may assume $k$ algebraically closed (\'etale base
change to $k^{sep}=\bar k$ does not change any of the ramification
filtrations involved).  By the principle of invariance of local invariants
under formal isomorphism (\Cref{section:ramification_after_blowup_is_tame},
discussion preceding the $r=1$ theorem) the ramification of $\cP^{[d]}$ at
$\eta_\ndls$ depends only on the completed local ring
$\widehat{\Oo}_{C,P}\cong k[[x]]$ together with the restriction of $\cP$ to
its punctured formal neighborhood, i.e.\ on the class
\[
[\wv a]\in W_r\big(k((x))\big)/\wp W_r\big(k((x))\big)
\]
corresponding to $\cP|_{\Spec\widehat K_P}$ under
\Cref{prop:classification} (equivalently, on the local character $\rho$ of
$G_{\widehat K_P}$).  By \Cref{lem:reduction} choose a \emph{reduced}
representative $\wv f$ of $[\wv a]$; since $\cP$ has ramification bounded
by $d$ at $P$, \Cref{thm:schmid-witt} gives
\begin{equation}\label{eq:bound}
\mu(\wv f)\;=\;\max_{0\le l\le r-1}p^{\,r-1-l}m_l\;\le\;d-1 .
\end{equation}

Now set $y:=x^{-1}$ and consider
\[
C\;=\;\bP^1_k,\qquad P\;=\;0,\qquad
U\;=\;\bP^1_k\setminus\{0\}\;=\;\Spec k[y]\;\cong\;\A^1_k .
\]
The components $f_l$ of $\wv f$ are polynomials in $y$ without constant
term, so $\wv f\in W_r(k[y])$ and
$\cQ:=\cP_{\wv f}\to U$ (\Cref{prop:etale}) is a $G$-torsor on $U$ whose
completed local datum at $0$ is $[\wv f]=[\wv a]$.  By the invariance
principle it therefore suffices to prove the theorem for
$(\bP^1,\,d\cdot 0,\,\cQ)$; \emph{we assume from now on that}
\[
\cP=\cP_{\wv f}\ \text{over}\ U=\Spec k[y],\qquad
\wv f\ \text{reduced},\qquad \mu(\wv f)\le d-1 .
\]

\begin{rem*}[the case $C=\bP^1$ needs no formal-invariance principle]
When $C=\bP^1$ to begin with, the replacement of $\cP$ by $\cP_{\wv f}$ is
an honest isomorphism of torsors and the invariance principle is not
needed.  Indeed the natural map
\[
W_r(k[y])/\wp W_r(k[y])\longrightarrow
W_r\big(k((x))\big)/\wp W_r\big(k((x))\big)
\]
is \emph{injective}: suppose $\wv g\in W_r(k[y])$ and
$\wv g=\wp(\wv h)$ with $\wv h\in W_r(k((x)))$.  At level $0$,
$h_0^p-h_0=g_0\in k[y]=k[x^{-1}]$; writing $h_0=h^-+c+h^+$
(pole part, constant, part in $\mm$) and using
$\wp(h_0)=\wp(h^-)+\wp(c)+\wp(h^+)$ (Frobenius is additive), the
$\mm$-part of $g_0$ is $\wp(h^+)=0$; if $h^+\neq0$ its lowest-degree term
survives in $h^{+p}-h^+$, so $h^+=0$ and $h_0\in k[y]$.  Inductively, at
level $n$, $h_n^p-h_n=g_n-\Delta_n(h_{<n})\in k[y]$
(\Cref{lem:wp-triangular}, with $h_{<n}\in k[y]$ by induction), and the
same argument gives $h_n\in k[y]$.  Hence $\wv g\in\wp W_r(k[y])$.
Consequently any two Witt vectors in $W_r(k[y])$ with the same class over
$k((x))$ define isomorphic torsors on $U$ (\Cref{prop:classification}).
\end{rem*}

\subsection{Coordinates on the symmetric power and the blowup}\label{subsec:coords}

Since $U=\Spec k[y]\cong\A^1$, we have
\[
U^{d}=\Spec k[y_1,\dots,y_d],\qquad
U^{(d)}=\Spec k[\eps_1,\dots,\eps_d],\qquad
\eps_j:=e_j(y_1,\dots,y_d),
\]
with $e_j$ the elementary symmetric polynomials.  On the other side, $x$ is
a coordinate at $P=0$, the point $\ndls=d\cdot 0\in C^{(d)}$ has the affine
neighborhood $(\A^1_x)^{(d)}=\Spec k[e_1,\dots,e_d]$ with
$e_j:=e_j(x_1,\dots,x_d)$ vanishing at $\ndls$, and (as computed in
\Cref{section:ramification_after_blowup_is_tame} for the blowup of affine
space at the origin) the local ring
$R:=\Oo_{X_\ndls,\eta_\ndls}$ is a discrete valuation ring with
\[
v_{\eta}(e_j)=1\ (1\le j\le d),\qquad
\kappa(\eta_\ndls)=k\Big(\tfrac{e_1}{e_d},\dots,\tfrac{e_{d-1}}{e_d}\Big),
\qquad
\widehat R=\kappa(\eta_\ndls)[[\,e_d\,]] .
\]
Since $y_i=1/x_i$, the two coordinate systems are related by
\begin{equation}\label{eq:eps-vs-e}
\eps_j\;=\;e_j(y_1,\dots,y_d)\;=\;\frac{e_{d-j}(x_1,\dots,x_d)}{e_d(x_1,\dots,x_d)}
\;=\;\frac{e_{d-j}}{e_d}
\qquad(1\le j\le d,\ e_0:=1).
\end{equation}
In particular:
\begin{equation}\label{eq:eps-in-kappa}
\eps_j\in\kappa(\eta_\ndls)\subset\widehat R
\quad\text{for }1\le j\le d-1,
\qquad\text{while}\qquad
\eps_d=\frac{1}{e_d}\ \text{has}\ v_\eta(\eps_d)=-1 .
\end{equation}
Both systems generate the same function field
$K(C^{(d)})=k(e_1,\dots,e_d)=k(\eps_1,\dots,\eps_d)$, and we write
$\widehat K_\eta=\kappa(\eta_\ndls)((e_d))=\Frac(\widehat R)$.

\subsection{The datum of \texorpdfstring{$\cP^{[d]}$}{P[d]} is the \texorpdfstring{$d$}{d}-fold Witt sum}

Recall (\Cref{cor:sym-power-quotient}) that
$\cP^{[d]}=\cP^{\times d}/\,\Xi$ with $\Xi=\KG\rtimes S_d$, where
$\KG=\ker\big(\Sigma\colon G^d\to G\big)$; for affine $U$ this quotient is
the spectrum of the ring of invariants.

\begin{theorem}[the symmetric-power datum]\label{thm:datum}
With $\cP=\cP_{\wv f}$ over $U=\Spec k[y]$ as in
\Cref{subsec:reductions}, there is an isomorphism of $G$-torsors over
$U^{(d)}=\Spec k[\eps_1,\dots,\eps_d]$
\[
\cP^{[d]}\;\cong\;\cP_{\wv F},
\qquad\text{where}\qquad
\wv F\;:=\;\Wsum_{i=1}^{d}\ \wv f(y_i)\;\in\;W_r\big(k[\eps_1,\dots,\eps_d]\big)
\]
is the $d$-fold Witt sum of the pullbacks of $\wv f$ along the $d$
projections.  (Its components
$F_j=S^{(d)}_j\big(\wv f(y_1),\dots,\wv f(y_d)\big)$ are symmetric
polynomials in $y_1,\dots,y_d$, hence indeed lie in
$k[\eps_1,\dots,\eps_d]$.)
\end{theorem}

\begin{proof}
Write $A^{\otimes d}=k[y_1,\dots,y_d]$.  The fiber product
$\cP^{\times d}=p_1^{-1}\cP\times_{U^d}\dots\times_{U^d}p_d^{-1}\cP$ is the
spectrum of
\[
T\;=\;A^{\otimes d}\big[\wv X_1,\dots,\wv X_d\big]\,/\,
\big(\wp(\wv X_i)-\wv f(y_i)\ :\ 1\le i\le d\big),
\]
a $G^d$-torsor over $U^d$, where $\wv g_i\in G$ acts by
$\wv X_i\mapsto\wv X_i+_W\wv g_i$, and $S_d$ acts by permuting
simultaneously the $y_i$ and the batches $\wv X_i$.

\emph{Step 1: $\KG$-invariants.}
Put $\wv Y:=\Wsum_{i=1}^{d}\wv X_i$.  Exactly as in the proof of
\Cref{lem:additivity} (using \Cref{lem:triangular}(1) for the batch
$\wv X_1$), the substitution
$X_{1,l}\mapsto Y_l$ is a triangular unipotent change of polynomial
generators, $A^{\otimes d}[\wv X_1,\dots,\wv X_d]
=A^{\otimes d}[\wv Y,\wv X_2,\dots,\wv X_d]$, and the same two-sided
ideal computation (with $\wv X_1=\wv Y-_W\wv X_2-_W\dots-_W\wv X_d$ and
additivity of $\wp$) identifies the defining ideal with
\[
\big(\wp(\wv Y)-\wv F,\ \ \wp(\wv X_j)-\wv f(y_j)\ :\ 2\le j\le d\big).
\]
Hence
\[
T\;\cong\;
\underbrace{A^{\otimes d}[\wv Y]/\big(\wp(\wv Y)-\wv F\big)}_{=:B}
\ \otimes_{A^{\otimes d}}\ T',
\qquad
T':=A^{\otimes d}[\wv X_2,\dots,\wv X_d]/\big(\wp(\wv X_j)-\wv f(y_j)\big),
\]
and $\Spec T'\to U^d$ is a $G^{d-1}$-torsor (\Cref{prop:etale}).
Parametrize $\KG\cong G^{d-1}$ by
$(\wv g_2,\dots,\wv g_d)\mapsto
\big(-\!\Wsum_{j\ge2}\wv g_j,\ \wv g_2,\dots,\wv g_d\big)$.
Under this action $\wv Y\mapsto\wv Y+_W\Wsum_i\wv g_i=\wv Y$ is invariant,
while $(\wv g_2,\dots,\wv g_d)$ acts on $T'$ through the full torsor
action of $G^{d-1}$.  By \Cref{lem:torsor-invariants},
\[
T^{\KG}\;=\;B\otimes_{A^{\otimes d}}(T')^{G^{d-1}}
\;=\;B\otimes_{A^{\otimes d}}A^{\otimes d}\;=\;B .
\]
The residual action of $G=G^d/\KG$ on $B$ is $\wv Y\mapsto\wv Y+_W\wv g$.

\emph{Step 2: $S_d$-invariants.}
The $S_d$-action on $T$ preserves the subring $B=T^{\KG}$ (it normalizes
$\KG$ in $\Xi$); explicitly it permutes the $y_i$ and fixes each component
$Y_l$ of $\wv Y$ (the Witt sum $\Wsum_i\wv X_i$ is invariant under
permuting its summands), and it fixes each component of $\wv F$ (a
symmetric polynomial).  By \Cref{prop:etale}(1), $B$ is a free
$A^{\otimes d}$-module with basis the monomials
$\prod_lY_l^{\,c_l}$ ($0\le c_l\le p-1$), and $S_d$ acts on $B$
coefficientwise with respect to this basis.  Hence
\[
B^{S_d}
=\big(A^{\otimes d}\big)^{S_d}\big[\wv Y\big]/\big(\wp(\wv Y)-\wv F\big)
=k[\eps_1,\dots,\eps_d]\big[\wv Y\big]/\big(\wp(\wv Y)-\wv F\big),
\]
which is precisely the ring of $\cP_{\wv F}$ over
$U^{(d)}=\Spec k[\eps_1,\dots,\eps_d]$.  Therefore
$\cP^{[d]}=\Spec T^{\Xi}=\Spec B^{S_d}=\cP_{\wv F}$, compatibly with the
$G$-actions.
\end{proof}

\begin{rem*}
For $r=1$ this is literally the computation of the thesis
(\Cref{section:ramification_after_blowup_is_tame}): $\wv Y$ is the single
variable $Y=X_1+\dots+X_d$ and $\wv F$ the power sum $\sum_if(y_i)$.  The
Witt sum is the correct $p^r$-analogue: the ``carries'' of Witt addition,
which obstructed the inductive approach, are now explicit polynomial terms
that the next two lemmas control.
\end{rem*}

% #############################################################################
\section{Integrality at the exceptional divisor, and the proof}\label{sec:proof}
% #############################################################################

\begin{lemma}[isobaric degree bound]\label{lem:isobaric-bound}
Let $\wv f$ be reduced with pole orders $m_l$ and weighted pole order
$\mu=\mu(\wv f)$, and let
$F_j=S^{(d)}_j\big(\wv f(y_1),\dots,\wv f(y_d)\big)$ be the components of
the Witt sum $\wv F$ of \Cref{thm:datum}.  Then $F_j$ is a symmetric
polynomial in $y_1,\dots,y_d$ of total degree
\[
\deg F_j\;\le\;p^{\,j-r+1}\mu\;\le\;\mu
\qquad(0\le j\le r-1).
\]
\end{lemma}

\begin{proof}
By \Cref{lem:isobaric}, every monomial of $S^{(d)}_j$ has the form
$\prod_{i,l}\big(X^{(i)}_l\big)^{n_{i,l}}$ with
$\sum_{i,l}n_{i,l}\,p^{\,l}=p^{\,j}$.  Substituting
$X^{(i)}_l\mapsto f_l(y_i)$, a polynomial of degree $m_l$ in $y_i$, the
resulting term has total degree
\[
\sum_{i,l}n_{i,l}\,m_l
\;\le\;\sum_{i,l}n_{i,l}\,p^{\,l}\cdot\max_{l}\frac{m_l}{p^{\,l}}
\;=\;p^{\,j}\cdot p^{\,1-r}\max_l\,p^{\,r-1-l}m_l
\;=\;p^{\,j-r+1}\mu\;\le\;\mu,
\]
using $j\le r-1$.  Symmetry of $F_j$ is clear since permuting the $y_i$
permutes the summands of the Witt sum.
\end{proof}

\begin{lemma}[low-degree symmetric polynomials are residue-field constants]\label{lem:integrality}
Let $F\in k[y_1,\dots,y_d]^{S_d}$ be a symmetric polynomial of total degree
$D<d$.  Then $F\in k[\eps_1,\dots,\eps_{d-1}]$ ($\eps_d$ does not occur),
and consequently, under \cref{eq:eps-vs-e,eq:eps-in-kappa}, the image of
$F$ in $\widehat K_\eta$ lies in the residue field
$\kappa(\eta_\ndls)\subset\widehat R$; in particular $v_\eta(F)\ge0$.
\end{lemma}

\begin{proof}
By the fundamental theorem of symmetric polynomials, $F$ is a polynomial
in $\eps_1,\dots,\eps_d$; we claim only $\eps_1,\dots,\eps_D$ occur.
Recall the standard algorithm: if $y^{\lambda}
=y_1^{\lambda_1}\cdots y_s^{\lambda_s}$
($\lambda_1\ge\dots\ge\lambda_s>0$) is the lexicographically leading
monomial of $F$, with coefficient $c$, one subtracts
$c\,\eps_{1}^{\lambda_1-\lambda_2}\eps_2^{\lambda_2-\lambda_3}\cdots
\eps_s^{\lambda_s}$, obtaining a symmetric polynomial of the same or
smaller degree with a lexicographically smaller leading term, and
iterates.  A monomial of degree $\le D$ involves at most $D$ distinct
variables, so $s\le D$ at every step, and only
$\eps_1,\dots,\eps_D$ with $D<d$ are ever used.  The last statement now
follows from \cref{eq:eps-in-kappa}: each $\eps_j$ with $j\le d-1$ maps to
the element $e_{d-j}/e_d$ of the coefficient field
$\kappa(\eta_\ndls)$ of $\widehat R=\kappa(\eta_\ndls)[[e_d]]$.
\end{proof}

\begin{rem*}
The hypothesis $D<d$ is exactly what excludes
$\eps_d=1/e_d$, the unique pole at $\eta_\ndls$ among the
$\eps_j$; this is the multi-level generalization of the thesis's $r=1$
lemma on $\sum_if(y_i)$, and it is sharp --- see the concluding remarks.
\end{rem*}

\subsubsection*{Proof of \Cref{theorem:ramification_after_blowup_is_tame_for_p_power}}

\begin{proof}
By the reductions of \Cref{subsec:reductions} we may assume
$k=\bar k$, $C=\bP^1$, $P=0$, $U=\Spec k[y]$ and $\cP=\cP_{\wv f}$ with
$\wv f$ reduced and $\mu(\wv f)\le d-1$ (this used
\Cref{lem:reduction} and the Schmid--Witt bound,
\Cref{thm:schmid-witt}).

By \Cref{thm:datum},
$\cP^{[d]}\cong\cP_{\wv F}$ over $U^{(d)}$ with
$\wv F=\Wsum_{i=1}^d\wv f(y_i)$.  Restricting to the completed field at
the generic point of the exceptional divisor
(\Cref{subsec:coords}), the torsor
$\cP^{[d]}|_{\Spec\widehat K_\eta}$ is $\cP_{\wv F'}$ where
$\wv F'\in W_r(\widehat K_\eta)$ is the image of $\wv F$.

By \Cref{lem:isobaric-bound}, each component $F_j$ is a symmetric
polynomial in $y_1,\dots,y_d$ of total degree $\le\mu(\wv f)\le d-1<d$;
hence by \Cref{lem:integrality} each $F'_j$ lies in
$\kappa(\eta_\ndls)\subset\widehat R$.  Therefore
\[
\wv F'\in W_r\big(\widehat R\big),
\]
and by \Cref{cor:integral-unramified} every field factor of
$\cP_{\wv F'}$ is unramified over $\widehat R$.  Since ramification at
$\eta_\ndls$ is computed exactly on this restriction
(\Cref{definition:tamely_ramified_in_codim_1} and the surrounding
discussion), $\cP^{[d]}$ is unramified at $\eta_\ndls$.
\end{proof}

% #############################################################################
\section*{Concluding remarks}
% #############################################################################

\begin{rem*}[sharpness, and the boundary examples]
\Cref{lem:isobaric-bound} is in fact an equality,
$\deg F_{r-1}=\mu(\wv f)$: the extremal monomial
$\big(X^{(i)}_{l_0}\big)^{p^{\,r-1-l_0}}$ realizing
$\mu=p^{\,r-1-l_0}m_{l_0}$ occurs in $S^{(d)}_{r-1}$ with a nonzero
coefficient mod $p$, and its top-degree part cannot cancel.  Hence the
hypothesis $\mu\le d-1$ is optimal:
if $\mu\ge d$ the top component of the datum genuinely involves
$\eps_d=1/e_d$ and acquires a pole at $\eta_\ndls$.  This matches the
boundary phenomena recorded in the thesis:
for $p=2$, $r=2$, $m_0=1$ (so $\mu=2$) the symmetric square ($d=2$) is
ramified at $\eta_\ndls$ (\Cref{ex:p2-counterexample}), while $d=3>\mu$ is
unramified; and the machine verifications in
\texttt{q\_ramification\_along\_VF/general\_witt\_proof}
(\texttt{witt\_general.py}, \texttt{witt\_r3.py}) confirm
$\deg F_j=p^{\,j-r+1}\mu$ exactly for $r=2,3$ and various $p$.
\end{rem*}

\begin{rem*}[why the induction stalled]
The inductive route factored $\cP\to\cP'\to U$ through the intermediate
curve and tried to reconstruct the ramification of $\cP^{[d]}$ from that
of $\cP'^{[d]}$ across the base change $q$; the loss of information there
(the ``absorption'' phenomenon of \texttt{a\_wrong\_path.tex}) is
precisely the Witt carry.  In the present proof the carry is not
estimated across a base change but computed in closed form inside the
symmetric power, where the isobaric grading shows it obeys the same
degree bound as the leading terms.  No induction on $r$, and no Swan
conductor theory, is needed.
\end{rem*}

\begin{rem*}[integration into the thesis]
\begin{enumerate}
  \item \Cref{sec:witt} and \Cref{sec:asw} belong in the Preliminaries,
        next to Kummer and Artin--Schreier theories
        (\Cref{theorem:kummer_ramification},
        \Cref{theorem:artin_schreier_ramification});
        \Cref{thm:schmid-witt} generalizes the latter and can replace it.
  \item \Cref{subsec:reductions,subsec:coords}, \Cref{thm:datum},
        \Cref{lem:isobaric-bound,lem:integrality} and the final proof
        replace the unfinished proof of
        \Cref{theorem:ramification_after_blowup_is_tame_for_p_power} in
        \texttt{RamificationAfterBlowupIsTame.tex}.  Note that
        \Cref{thm:datum} at $r=1$ also subsumes the Artin--Schreier half
        of the theorem preceding
        \Cref{cor:ramification_after_blowup_is_tame}, so the two proofs
        can be merged if desired.
  \item The deduction of the general abelian case
        (\Cref{theorem:torsors_after_blowup_is_tame}) is unchanged:
        decompose $G$ by the structure theorem and combine the cyclic
        cases via \Cref{prop:quotient_torsors} and
        \Cref{lemma:bounding_ramification_of_contracted_product} (or
        \Cref{prop:ramification-external-product}), exactly as outlined
        in the thesis.
\end{enumerate}
\end{rem*}

\printbibliography

\end{document}
